% Данный файл распространяетсяя под лицензией CC BY 4.0. Текст лицензии размещён на https://creativecommons.org/licenses/by/4.0/
% (c) Кафедра системного программирования, 2025

\documentclass[a4paper]{article}

\usepackage[a4paper, top=8mm, bottom=8mm, left=8mm, right=8mm]{geometry}

\usepackage{polyglossia}
\setdefaultlanguage[babelshorthands=true]{russian}
\setotherlanguage{english}

\usepackage{fontspec}
\setmainfont{FreeSerif}
\newfontfamily{\russianfonttt}[Scale=0.7]{DejaVuSansMono}

\usepackage[tiny, compact]{titlesec}

\usepackage{titling}
\setlength{\droptitle}{-1cm}
\pretitle{\begin{center}\begin{bfseries}\Large}
\posttitle{\par\end{bfseries}\end{center}}
\preauthor{\begin{center}\normalsize}
\postauthor{\par\end{center}\vspace{-1.8cm}}

\usepackage{hyperref}
\usepackage{bookmark}
\usepackage{csquotes}
\usepackage{booktabs}
\usepackage{longtable}
\usepackage{multirow}

\title{Реализация алгоритмов поиска кратчайшего пути в графе на основе GraphBLAS API}

\author{Шаманаев Иван Сергеевич}

\date{}

\begin{document}

\maketitle

\begin{flushright}
    Группа: \emph{2025.Б72-мм}

    Кафедра: \emph{кафедра системного программирования}

    Научный руководитель: \emph{Григорьев Семён Вячеславович,}\\
    \emph{доцент кафедры системного программирования,}\\
    \emph{кандидат физико-математических наук}

    Номер семестра практики: \emph{2}
\end{flushright}


\section{Постановка задачи}
Работа выполнялась в рамках учебной ознакомительной практики и направлена на знакомство с GraphBLAS API (стандарт API для графовых операций через линейную алгебру) и сравнительный анализ производительности реализаций алгоритмов поиска кратчайшего пути на языке C.

\textbf{Целью работы является} реализация алгоритмов поиска кратчайшего пути в графе на основе GraphBLAS API, экспериментальное сравнение собственных реализаций между собой, а также с готовой реализацией из библиотеки LAGraph.

Для достижения поставленной цели были решены следующие задачи:
\begin{itemize}
    \item выполнить обзор подходов к интеграции GraphBLAS API на языке C;
    \item реализовать алгоритмы поиска кратчайшего пути с использованием SuiteSparse:GraphBLAS;
    \item провести замеры производительности разработанных решений и готовых реализаций из LAGraph.
\end{itemize}
\section{Описание предлагаемого решения}

Проект представляет собой сравнительное тестирование для анализа трёх алгоритмов SSSP (Single-Source Shortest Path~--- поиск кратчайшего пути из одной вершины), реализованных с использованием библиотек SuiteSparse:GraphBLAS и LAGraph на языке C.\footnote{URL: \url{https://github.com/DrTimothyAldenDavis/GraphBLAS} (дата обращения: 30.07.2026).}\footnote{URL: \url{https://github.com/DrTimothyAldenDavis/LAGraph} (дата обращения: 30.07.2026).}

Алгоритм \textbf{Delta-Stepping} реализован через готовую функцию библиотеки LAGraph
\texttt{LAGr\_SingleSourceShortestPath} с параметром $\delta$, определяющим ширину
«шага» для группировки вершин. Не поддерживает отрицательные веса рёбер.

Алгоритм \textbf{Algebraic Bellman-Ford} реализован на чистом GraphBLAS с
использованием min-plus полукольца: моноид \texttt{GrB\_MIN\_FP64} и операция
\texttt{GrB\_PLUS\_FP64}. За $n-1$ итерацию выполняется умножение вектора расстояний
на матрицу смежности через \texttt{GrB\_vxm}. Поддерживает отрицательные веса
(при отсутствии отрицательных циклов).

Алгоритм \textbf{Дейкстры} реализован на чистом GraphBLAS с самописной приоритетной
очередью (бинарная куча). На каждом шаге выбирается вершина с минимальным расстоянием,
после чего выполняется релаксация смежных рёбер через \texttt{GrB\_vxm}.
Только для неотрицательных весов.

\section{Эксперименты}

Для сравнения производительности реализованных алгоритмов были проведены следующие испытания.

\textbf{Среда выполнения.} 
Сборка и тестирование выполнялись на ноутбуке Asus ROG Flow x16 с процессором Intel Core i9-13900H и 32~ГБ оперативной памяти частотой 4800~МГц. 
При замерах для более стабильного результата выполнялось по 40 запусков для каждого графа; первые 10 считались прогревочными и не учитывались в результатах. Во время замеров был включён производительный режим, фоновые приложения отключены. Питание ноутбука было от сети с полным зарядом.
Окружения, в которых производилось тестирование производительности:
\begin{itemize}
    \item Fedora 43;
    \item MSYS2 / MinGW64 под Windows.
\end{itemize}


Для каждой платформы требуется установка SuiteSparse GraphBLAS 10.3.1 и LAGraph 1.2.1.

Характеристики графов, использованных для тестирования, приведены в
таблице~\ref{tab:datasets}. Все графы взяты из коллекции SuiteSparse Matrix
Collection.\footnote{URL: \url{https://sparse.tamu.edu} (дата обращения: 30.07.2026).}

\begin{table}[h!]
    \centering
    \caption{Используемые наборы данных}
    \label{tab:datasets}
    \begin{tabular}{lrr}
        \toprule
        \textbf{Граф} & \textbf{Вершины} & \textbf{Рёбра} \\
        \midrule
        amazon0601   & 403\,394  & 3\,387\,388  \\
        com-Amazon   & 334\,863  &   925\,872  \\
        com-Orkut    & 3\,072\,441 & 117\,185\,083 \\
        NLR          & 4\,163\,763 & 12\,487\,976 \\
        roadNet-CA   & 1\,971\,281 & 2\,766\,607  \\
        web-Google   & 916\,428  & 5\,105\,039  \\
        wiki-topcats & 1\,791\,489 & 28\,511\,807 \\
        sinc18       & 16\,428   & 973\,826     \\
        \bottomrule
    \end{tabular}
\end{table}

\textbf{Сравнительное тестирование.} Основная программа (\texttt{bin/sssp\_benchmark}) последовательно
запускает все три алгоритма на загруженном графе, замеряя время выполнения каждого
через \texttt{LAGraph\_WallClockTime()}. Для Delta-Stepping варьировался параметр $\delta$
(0.5, 1, 2, 3, 5, 10). Диапазон выбран так, чтобы покрыть малые значения
(при которых алгоритм приближается к обычному Bellman-Ford) и большие
(при которых увеличиваются размеры «шагов», но снижается точность группировки).
Значения, кратные 1, удобны для интерпретации; 0.5 добавлено для проверки
поведения при дробном шаге; 10 выбрано как верхняя граница, после которой
выигрыш в производительности на тестовых графах перестаёт расти.
Для каждого графа и каждого алгоритма выполнялось 30 замеров, результаты
усреднялись. Для каждого набора замеров вычислялось стандартное отклонение
(Std), показывающее разброс времени выполнения относительно среднего.
В столбце «Время (мс)» после $\pm$ указан 95~\% доверительный интервал
среднего (95~\% CI).

Для графов с отрицательными весами рёбер (sinc18) Delta-Stepping (LAGraph) и Dijkstra
автоматически пропускаются, выполняется только Algebraic Bellman-Ford,
поддерживающий такие графы.

\textbf{Валидация результатов.} После выполнения всех алгоритмов производится
попарное сравнение векторов расстояний с погрешностью $10^{-6}$
(\texttt{VALIDATOR\_EPSILON}). Если результаты расходятся, выводится
предупреждение. Для графа без отрицательных весов все три алгоритма должны
давать одинаковые расстояния.

Результаты замеров на Linux (Fedora~43) представлены в таблице~\ref{tab:bench-linux},
на Windows~(MSYS2 / MinGW64)~--- в таблице~\ref{tab:bench-win}.


\begin{table}[h!]
\centering
\begin{minipage}{0.48\textwidth}
\centering
\footnotesize
\caption{Результаты тестирования (Linux, Fedora~43)}
\label{tab:bench-linux}
\begin{tabular}{llcc}
\toprule
\textbf{Граф} & \textbf{Конфигурация} & \textbf{Время (мс)} & \textbf{Std (\%)} \\
\midrule
\multirow{8}{*}{amazon0601}
    & Alg. BF & $58.40 \pm 0.41$ & 1.90 \\
    & Dijkstra & $1383.84 \pm 2.66$ & 0.52 \\
    & DS $\delta=0.5$ & $37.09 \pm 0.20$ & 1.46 \\
    & DS $\delta=1$ & $37.74 \pm 0.30$ & 2.12 \\
    & DS $\delta=2$ & $48.00 \pm 0.31$ & 1.73 \\
    & DS $\delta=3$ & $54.79 \pm 0.53$ & 2.57 \\
    & DS $\delta=5$ & $53.87 \pm 0.45$ & 2.23 \\
    & DS $\delta=10$ & $51.68 \pm 0.29$ & 1.51 \\
\midrule
\multirow{8}{*}{com-Amazon}
    & Alg. BF & $48.21 \pm 0.46$ & 2.57 \\
    & Dijkstra & $886.93 \pm 10.83$ & 3.27 \\
    & DS $\delta=0.5$ & $33.97 \pm 0.37$ & 2.94 \\
    & DS $\delta=1$ & $37.60 \pm 0.31$ & 2.23 \\
    & DS $\delta=2$ & $43.55 \pm 0.29$ & 1.81 \\
    & DS $\delta=3$ & $47.32 \pm 0.24$ & 1.37 \\
    & DS $\delta=5$ & $49.57 \pm 0.51$ & 2.74 \\
    & DS $\delta=10$ & $47.36 \pm 0.31$ & 1.77 \\
\midrule
\multirow{8}{*}{com-Orkut}
    & Alg. BF & $798.47 \pm 1.36$ & 0.46 \\
    & Dijkstra & $20240.79 \pm 163.35$ & 2.16 \\
    & DS $\delta=0.5$ & $125.05 \pm 0.57$ & 1.23 \\
    & DS $\delta=1$ & $128.57 \pm 0.92$ & 1.91 \\
    & DS $\delta=2$ & $137.64 \pm 1.12$ & 2.18 \\
    & DS $\delta=3$ & $140.85 \pm 0.85$ & 1.62 \\
    & DS $\delta=5$ & $113.95 \pm 0.88$ & 2.08 \\
    & DS $\delta=10$ & $95.62 \pm 1.10$ & 3.07 \\
\midrule
\multirow{8}{*}{NLR}
    & Alg. BF & $14217.48 \pm 9.95$ & 0.19 \\
    & Dijkstra & $22731.97 \pm 144.84$ & 1.71 \\
    & DS $\delta=0.5$ & $1315.01 \pm 9.42$ & 1.92 \\
    & DS $\delta=1$ & $1563.72 \pm 4.75$ & 0.81 \\
    & DS $\delta=2$ & $1455.72 \pm 6.63$ & 1.22 \\
    & DS $\delta=3$ & $1410.03 \pm 2.08$ & 0.40 \\
    & DS $\delta=5$ & $1360.90 \pm 1.71$ & 0.34 \\
    & DS $\delta=10$ & $1325.95 \pm 1.24$ & 0.25 \\
\midrule
\multirow{8}{*}{roadNet-CA}
    & Alg. BF & $277.17 \pm 1.41$ & 1.36 \\
    & Dijkstra & $3894.41 \pm 39.73$ & 2.73 \\
    & DS $\delta=0.5$ & $50.53 \pm 0.73$ & 3.86 \\
    & DS $\delta=1$ & $47.32 \pm 0.44$ & 2.49 \\
    & DS $\delta=2$ & $43.33 \pm 0.72$ & 4.43 \\
    & DS $\delta=3$ & $41.48 \pm 0.52$ & 3.38 \\
    & DS $\delta=5$ & $38.69 \pm 0.43$ & 3.00 \\
    & DS $\delta=10$ & $35.74 \pm 0.24$ & 1.79 \\
\midrule
\multirow{8}{*}{web-Google}
    & Alg. BF & $108.28 \pm 0.74$ & 1.83 \\
    & Dijkstra & $2616.09 \pm 22.83$ & 2.34 \\
    & DS $\delta=0.5$ & $49.21 \pm 0.50$ & 2.74 \\
    & DS $\delta=1$ & $51.10 \pm 0.48$ & 2.50 \\
    & DS $\delta=2$ & $65.87 \pm 0.56$ & 2.29 \\
    & DS $\delta=3$ & $73.14 \pm 0.47$ & 1.72 \\
    & DS $\delta=5$ & $75.51 \pm 0.44$ & 1.55 \\
    & DS $\delta=10$ & $72.88 \pm 0.41$ & 1.50 \\
\midrule
\multirow{8}{*}{wiki-topcats}
    & Alg. BF & $3194.51 \pm 11.61$ & 0.97 \\
    & Dijkstra & $8674.19 \pm 73.08$ & 2.26 \\
    & DS $\delta=0.5$ & $114.34 \pm 0.69$ & 1.62 \\
    & DS $\delta=1$ & $95.07 \pm 0.73$ & 2.06 \\
    & DS $\delta=2$ & $95.20 \pm 0.80$ & 2.25 \\
    & DS $\delta=3$ & $98.49 \pm 0.58$ & 1.57 \\
    & DS $\delta=5$ & $91.69 \pm 0.70$ & 2.05 \\
    & DS $\delta=10$ & $63.73 \pm 1.03$ & 4.33 \\
\midrule
\multirow{1}{*}{sinc18$^*$}
    & Alg. BF & $3495.55 \pm 34.33$ & 2.63 \\
\bottomrule
\end{tabular}
\smallskip

\noindent\emph{$^*$~--- sinc18 содержит отрицательные веса, поэтому Dijkstra и Delta-Stepping не запускались.}
\end{minipage}
\hfill
\begin{minipage}{0.48\textwidth}
\centering
\footnotesize
\caption{Результаты тестирования (Windows, MSYS2 / MinGW64)}
\label{tab:bench-win}
\begin{tabular}{llcc}
\toprule
\textbf{Граф} & \textbf{Конфигурация} & \textbf{Время (мс)} & \textbf{Std (\%)} \\
\midrule
\multirow{8}{*}{amazon0601}
    & Alg. BF & $221.67 \pm 14.39$ & 17.39 \\
    & Dijkstra & $7763.40 \pm 15.88$ & 0.55 \\
    & DS $\delta=0.5$ & $86.80 \pm 3.10$ & 9.56 \\
    & DS $\delta=1$ & $85.80 \pm 2.11$ & 6.60 \\
    & DS $\delta=2$ & $91.90 \pm 2.36$ & 6.89 \\
    & DS $\delta=3$ & $93.23 \pm 2.29$ & 6.58 \\
    & DS $\delta=5$ & $92.77 \pm 2.74$ & 7.92 \\
    & DS $\delta=10$ & $83.93 \pm 1.34$ & 4.29 \\
\midrule
\multirow{8}{*}{com-Amazon}
    & Alg. BF & $134.33 \pm 0.92$ & 1.84 \\
    & Dijkstra & $4920.80 \pm 30.95$ & 1.68 \\
    & DS $\delta=0.5$ & $74.63 \pm 2.21$ & 7.95 \\
    & DS $\delta=1$ & $76.40 \pm 2.12$ & 7.45 \\
    & DS $\delta=2$ & $75.30 \pm 1.88$ & 6.68 \\
    & DS $\delta=3$ & $78.80 \pm 2.70$ & 9.16 \\
    & DS $\delta=5$ & $76.67 \pm 1.28$ & 4.47 \\
    & DS $\delta=10$ & $71.70 \pm 2.27$ & 8.49 \\
\midrule
\multirow{8}{*}{com-Orkut}
    & Alg. BF & $966.30 \pm 5.14$ & 1.43 \\
    & Dijkstra & $70873.00 \pm 1494.83$ & 5.65 \\
    & DS $\delta=0.5$ & $224.77 \pm 2.08$ & 2.48 \\
    & DS $\delta=1$ & $233.83 \pm 4.19$ & 4.80 \\
    & DS $\delta=2$ & $232.03 \pm 1.84$ & 2.12 \\
    & DS $\delta=3$ & $231.13 \pm 1.16$ & 1.35 \\
    & DS $\delta=5$ & $217.03 \pm 3.29$ & 4.05 \\
    & DS $\delta=10$ & $178.23 \pm 1.84$ & 2.77 \\
\midrule
\multirow{8}{*}{NLR}
    & Alg. BF & $30267.10 \pm 108.68$ & 0.96 \\
    & Dijkstra & $82856.93 \pm 319.52$ & 1.03 \\
    & DS $\delta=0.5$ & $1886.17 \pm 19.87$ & 2.82 \\
    & DS $\delta=1$ & $2017.40 \pm 28.14$ & 3.74 \\
    & DS $\delta=2$ & $1714.60 \pm 17.81$ & 2.78 \\
    & DS $\delta=3$ & $1623.07 \pm 13.16$ & 2.17 \\
    & DS $\delta=5$ & $1522.53 \pm 12.50$ & 2.20 \\
    & DS $\delta=10$ & $1488.80 \pm 12.80$ & 2.30 \\
\midrule
\multirow{8}{*}{roadNet-CA}
    & Alg. BF & $966.47 \pm 7.76$ & 2.15 \\
    & Dijkstra & $7377.90 \pm 81.33$ & 2.95 \\
    & DS $\delta=0.5$ & $135.17 \pm 0.68$ & 1.35 \\
    & DS $\delta=1$ & $112.47 \pm 1.00$ & 2.38 \\
    & DS $\delta=2$ & $92.63 \pm 1.92$ & 5.56 \\
    & DS $\delta=3$ & $87.73 \pm 2.37$ & 7.24 \\
    & DS $\delta=5$ & $71.50 \pm 1.09$ & 4.10 \\
    & DS $\delta=10$ & $64.80 \pm 1.35$ & 5.57 \\
\midrule
\multirow{8}{*}{web-Google}
    & Alg. BF & $277.60 \pm 4.99$ & 4.82 \\
    & Dijkstra & $11998.30 \pm 22.37$ & 0.50 \\
    & DS $\delta=0.5$ & $120.97 \pm 1.59$ & 3.51 \\
    & DS $\delta=1$ & $121.93 \pm 1.59$ & 3.50 \\
    & DS $\delta=2$ & $129.07 \pm 1.15$ & 2.38 \\
    & DS $\delta=3$ & $131.27 \pm 0.98$ & 2.00 \\
    & DS $\delta=5$ & $130.07 \pm 1.91$ & 3.94 \\
    & DS $\delta=10$ & $120.77 \pm 3.80$ & 8.44 \\
\midrule
\multirow{8}{*}{wiki-topcats}
    & Alg. BF & $7509.20 \pm 81.29$ & 2.90 \\
    & Dijkstra & $36262.53 \pm 102.93$ & 0.76 \\
    & DS $\delta=0.5$ & $299.27 \pm 1.34$ & 1.20 \\
    & DS $\delta=1$ & $265.47 \pm 3.26$ & 3.29 \\
    & DS $\delta=2$ & $250.80 \pm 7.24$ & 7.73 \\
    & DS $\delta=3$ & $250.97 \pm 1.76$ & 1.88 \\
    & DS $\delta=5$ & $227.73 \pm 2.57$ & 3.02 \\
    & DS $\delta=10$ & $169.43 \pm 1.56$ & 2.47 \\
\midrule
\multirow{1}{*}{sinc18$^*$}
    & Alg. BF & $7845.23 \pm 246.55$ & 8.42 \\
\bottomrule
\end{tabular}
\smallskip

\noindent\emph{$^*$~--- sinc18 содержит отрицательные веса, поэтому Dijkstra и Delta-Stepping не запускались.}
\end{minipage}
\end{table}

Из результатов видно, что Delta-Stepping (DS) стабильно превосходит Algebraic Bellman-Ford
на всех графах, особенно на крупных (NLR~--- в ~10 раз, wiki-topcats~--- в ~30--50 раз).
Dijkstra на чистом GraphBLAS с самописной очередью значительно уступает обоим
алгоритмам, что связано с накладными расходами на извлечение элементов из
GraphBLAS-векторов на каждой итерации. Оптимальное значение $\delta$ зависит от
структуры графа: для разреженных графов (roadNet-CA) лучше большие $\delta$,
для плотных (com-Orkut)~--- малые.

\section{Заключение}

В ходе выполнения работы получены следующие результаты:

\begin{itemize}
    \item выполнен обзор подходов к интеграции GraphBLAS API на языке C, выбрана библиотека SuiteSparse:GraphBLAS;
    \item реализованы алгоритмы поиска кратчайшего пути (Algebraic Bellman-Ford, Дейкстра) с использованием SuiteSparse:GraphBLAS;
    \item проведены замеры производительности разработанных решений и готовых реализаций из LAGraph: Algebraic Bellman-Ford в большинстве случаев уступает по производительности Delta-Stepping, реализованному в LAGraph, и сильно превосходит алгоритм Дейкстры.
\end{itemize}

Репозиторий с реализацией: \url{https://github.com/IvRevver/graph_blas_sssp}.

Скрипты для подготовки окружения и запуска замеров, инструкции по сборке и воспроизведению результатов представлены в файле \texttt{README.md} репозитория.

\end{document}